
\begingroup

\begin{definition}[Regular scheme]
\label{def:regular-scheme}
\lean{DeJongAlterations.IsRegularScheme}
\leanok
On a locally Noetherian scheme, regularity means that every local ring is
regular. The stalkwise condition also makes sense without local
Noetherianity; all varieties in the alteration theorem are Noetherian.
\end{definition}

\endgroup

\begingroup

\begin{definition}[Strict normal crossings]
\label{def:snc}
\source{dejong96-alterations:page-0006}
\uses{def:regular-scheme}
\lean{DeJongAlterations.IsStrictNormalCrossings}
\leanok
On a Noetherian scheme $S$, a divisor $D$ has strict normal crossings if it
is reduced, $S$ is regular at all points of $D$, and the scheme-theoretic
intersection of every nonempty finite family of distinct irreducible components
is regular of codimension equal to the number of components. Each component
carries its reduced induced structure; their intersection is defined by the
sum of the component ideal sheaves, without taking its radical. Noetherianity
makes the set of components finite. Effective Cartier divisors are expressed
locally by a single non-zero-divisor, and pure codimension is measured by
ambient local-ring dimensions at the generic points of every intersection
component.
\end{definition}

\endgroup

\begingroup

\begin{definition}[Projectivity]
\label{def:projective-over}
\lean{DeJongAlterations.IsProjectiveOver}
\leanok
A morphism $p:X\to\operatorname{Spec}k$ is projective if it is the composite
of a closed immersion $X\hookrightarrow\mathbb P^n_k$ and the structure
morphism of $\mathbb P^n_k$, for some finite $n$.
\end{definition}

\endgroup

\begingroup

\begin{definition}[Alteration with boundary]
\label{def:alteration-output}
\uses{def:projective-over,def:regular-scheme,def:snc}
\lean{DeJongAlterations.AlterationWithBoundary}
\leanok
For a $k$-variety $X$ and a subset $Z$, an alteration with boundary consists
of $k$-varieties $Y,\overline Y$, a $k$-alteration $f:Y\to X$, a $k$-open
immersion $j:Y\hookrightarrow\overline Y$, and a reduced strict normal
crossings divisor $D$ on the regular projective variety $\overline Y$, with
support exactly $j(f^{-1}(Z))\cup(\overline Y\setminus j(Y))$.
\end{definition}

\endgroup

\begingroup

\begin{lemma}[Open subschemes]
\label{lem:regular-open}
\uses{def:regular-scheme}
\lean{DeJongAlterations.IsRegularScheme.of_isOpenImmersion}
\leanok
An open subscheme of a scheme with regular local rings has regular local rings.
\end{lemma}
\begin{proof}
An open immersion induces an isomorphism from the local ring of the target
at the image of a point to the local ring of the source at that point.
Regularity of a local ring is invariant under ring isomorphisms.
\end{proof}

\endgroup

\begingroup

\begin{lemma}[Actual divisor components and their finiteness]
\label{lem:divisor-components}
\uses{def:snc}
\lean{DeJongAlterations.divisorComponentOf,DeJongAlterations.divisorComponent_finite}
\leanok
Every irreducible component of the closed support of $D$ gives an ambient
closed component in the component type used above. If the ambient space is
Noetherian, this type is finite.
\end{lemma}
\begin{proof}
The inclusion of the closed support is a closed embedding. It takes each
closed irreducible component to a closed subset of the ambient space, and
the defining witness records its origin. A subspace of a Noetherian space
is Noetherian, so the support has finitely many irreducible components.
Their map to the component type is surjective by its defining witness;
thus that type is finite as well. In particular every subset of its
components can be represented by a finite set.
\end{proof}

\endgroup

\begingroup

\begin{corollary}[Dimension under an actual dense open immersion]
\label{cor:function-field-dimension-dense-open}
\source{dejong96-alterations:page-0019}
\lean{DeJongAlterations.Geometry.FunctionFieldDimension.topologicalKrullDim_eq_of_openImmersion,DeJongAlterations.Geometry.FunctionFieldDimension.topologicalKrullDim_le_one_of_dense_open}
\leanok
Let $Y$ be an integral scheme locally of finite type over $k$ and
$j:X\to Y$ a dominant open immersion. Then $\dim X=\dim Y$.
In particular $\dim X\leq1$ implies $\dim Y\leq1$.
\end{corollary}
\begin{proof}
The actual open image is nonempty by dominance. Choose a nonempty
affine neighborhood $U$ contained in that image. The scheme dimension
theorem shows that $Y$ and every nonempty affine open of $Y$, in
particular $U$, have the same dimension: both are the finite degree of
the same canonical function field. The universal property of the open
immersion lifts $U\to Y$ to an open immersion $U\to X$. Hence
$\dim Y=\dim U\leq\dim X$. The original open immersion also gives
$\dim X\leq\dim Y$, proving equality.
\end{proof}

\endgroup

\begingroup

\begin{lemma}[Alterations preserve dimension]
\label{lem:sm-dj-dim-alteration}
\source{dejong96-alterations:page-0012}
\uses{cor:function-field-dimension-dense-open}
\lean{DeJongAlterations.StableModification.DeJong.IsAlteration.toSurjective,DeJongAlterations.StableModification.DeJong.topologicalKrullDim_eq_of_isAlteration,DeJongAlterations.StableModification.DeJong.topologicalKrullDim_eq_of_isModification}
\leanok
Let $S$ be a scheme locally of finite type over a field and $f:X\to S$ an
alteration. Then $\dim X=\dim S$. In particular modifications preserve the
dimension.
\end{lemma}
\begin{proof}
The morphism $f$ is proper and dominant, hence surjective. Choose a nonempty
open $U\subseteq S$ such that $f^{-1}(U)\to U$ is finite. Then $f^{-1}(U)$ is
nonempty by surjectivity, and $f^{-1}(U)$ and $U$ are integral, being nonempty
open subschemes of integral schemes. The restriction $f^{-1}(U)\to U$ is finite
and surjective, hence dominant, so its function field extension is finite and
$\dim f^{-1}(U)=\dim U$, both dimensions being the transcendence degree of the
function field over $k$. A nonempty open subscheme of an integral scheme of
finite type over a field has the dimension of the whole scheme; applying this to
$f^{-1}(U)\subseteq X$ ($X$ is of finite type over $k$ through $f$) and to
$U\subseteq S$ gives $\dim X=\dim f^{-1}(U)=\dim U=\dim S$.
\end{proof}

\endgroup

\begingroup

\begin{lemma}[Varieties have finite dimension]
\label{lem:sm-dj-dim-finite}
\uses{def:variety}
\lean{DeJongAlterations.StableModification.DeJong.exists_nat_topologicalKrullDim_eq}
\leanok
The dimension of a variety over a field is a natural number.
\end{lemma}
\begin{proof}
It equals the transcendence degree of its function field over the ground
field, which is finite because the function field is finitely generated.
\end{proof}

\endgroup

\begingroup

\begin{lemma}[The actual dimension of a Cartier support component]
\label{lem:actual-cartier-component-dimension}
\lean{DeJongAlterations.Geometry.CartierBoundaryComponentDimension.component_affine_height,DeJongAlterations.Geometry.CartierBoundaryComponentDimension.component_dimension_add_one}
\leanok
Let $X$ be a variety over any field and let $D$ be an actual effective
Cartier ideal. For every native support component with actual reduced
vanishing-ideal closed subscheme $C$, one has $\dim C+1=\dim X$.
\end{lemma}
\begin{proof}
Choose a nonempty affine open meeting the component and write $I$ for
the component's actual vanishing ideal in its coordinate ring $A$.
The inverse image of the closed irreducible component is irreducible;
the ideal is radical, so $I$ is prime. The generic point of $V(I)$
maps to the generic point of the original support component. Actual
Cartier pure codimension gives ambient local dimension one there.
The native affine open stalk isomorphism and prime-height formula
give $\operatorname{ht}I=1$. The preceding exact affine dimension
formula yields $\dim(A/I)+1=\dim A$. The spectrum of this literal
quotient is an actual nonempty open chart of the actual reduced
component subscheme. Nonempty affine opens have the original global
finite-type variety dimensions, which gives the stated equality.
\end{proof}

\endgroup

\begingroup

\begin{theorem}[Dominant integral pullback preserves Cartier divisors]
\label{thm:cartier-dominant-pullback}
\source{dejong96-alterations:page-0018}
\lean{DeJongAlterations.Geometry.CartierDominantPullback.isEffectiveCartierDivisor_comap}
\leanok
Let $f:X\to Y$ be a dominant morphism of integral schemes.
If $D$ is an effective Cartier divisor on $Y$, then its actual
scheme-theoretic pullback is an effective Cartier divisor on $X$.
No finite-type, Noetherian, properness or flatness hypothesis is needed.
\end{theorem}
\begin{proof}
At $f(x)$ choose an affine Cartier neighborhood $V$ with full ideal
$(t)$, where $t$ is a non-zero-divisor. Choose a nonempty affine
neighborhood $U$ of $x$ inside $f^{-1}(V)$. Both section rings are
domains, and the preceding lemma shows that $t$ has nonzero image
in $\Gamma(U)$. That image is a non-zero-divisor. The exact affine
pullback lemma identifies the full pulled-back ideal on
$\operatorname{Spec}\Gamma(U)$ with the principal ideal generated
by this image. Its Cartier equation transports through the actual
open immersion to an affine neighborhood of $x$ in $X$. This proves
the required local statement at every point.
\end{proof}

\endgroup

\begingroup

\begin{theorem}[An actual regular codimension-one component is Cartier]
\label{thm:snc-regular-component-cartier}
\lean{DeJongAlterations.Geometry.SNCSubdivisorComponent.isRegularRing_of_spec,DeJongAlterations.Geometry.SNCSubdivisorComponent.isRegularRing_quotient_of_specIdealSheaf,DeJongAlterations.Geometry.SNCSubdivisorComponent.isEffectiveCartierDivisor_vanishingIdeal,DeJongAlterations.Geometry.SNCSubdivisorComponent.component_le_support,DeJongAlterations.Geometry.SNCSubdivisorComponent.component_isEffectiveCartierDivisor}
\leanok
Let $X$ be locally Noetherian, and let $C\subset X$ be irreducible
closed with pure codimension one. Suppose $X$ is regular along $C$
and the actual reduced closed subscheme of $C$ is regular.
Then its actual vanishing ideal is effective Cartier.
\end{theorem}
\begin{proof}
At a point of $C$, the actual closed-immersion stalk map is a
surjection between the two regular local rings. Its kernel is
the extension of the actual affine prime defining $C$, which
has height one by the affine-height lemma, and is therefore
nonzero in the domain given by the ambient regular local ring.
The kernel lemma constructs its regular principal equation.
Spread this equation by the preceding lemma. At a point outside
$C$, the restricted ideal is the unit ideal and has equation one.
These actual affine equations prove the Cartier assertion.
\leanok
\end{proof}

\endgroup

\begingroup

\begin{theorem}[Dominance from an integral scheme over a Dedekind domain forces flatness]
\label{thm:proper-flattening-dedekind-native-flat}
\lean{DeJongAlterations.Geometry.ProperFlatteningDedekind.flat_of_dominant_to_spec}
\leanok
Let $R$ be a Dedekind domain and $f:X\to\operatorname{Spec}R$ the
original dominant morphism from an integral scheme.  Then this same
morphism $f$ is flat.  No finite-type, properness, or relative-dimension
hypothesis on $f$ is required.
\end{theorem}
\begin{proof}
Cover $X$ by its actual nonempty affine neighborhoods.  Every such open
is integral, and its inclusion has nonempty open image in the irreducible
scheme $X$, hence is dominant.  Its composition with the original $f$
is therefore dominant.  Writing that actual affine map as
$\operatorname{Spec}B\to\operatorname{Spec}R$, both rings are domains.
The native dense-spectrum criterion says that its coefficient-map kernel
is contained in the nilradical of $R$, which is zero.  Thus the actual
map $R\to B$ is injective.  Consequently the original $R$-module $B$
is torsion free: a nonzero base coefficient maps to a nonzero element
of the domain $B$ and cannot annihilate a nonzero element of $B$.
The torsion-free flatness theorem over a Dedekind domain makes this very
$B$ flat over $R$.  The native affine criterion makes the actual affine
restriction flat.  Source locality on the chosen actual affine cover
proves native flatness of the original $f$.
\end{proof}

\endgroup

\begingroup

\begin{theorem}[Unconditional actual alteration over an arbitrary field]
\label{thm:sm-bridge-alteration-arbitrary-unconditional}
\lean{DeJongAlterations.StableModificationBridge.AlterationUnconditional.deJongAlterationTheorem}
\leanok
For every field $k$, with no perfectness or geometric integrality
assumption, the native actual \texttt{DeJongAlterationTheorem k}
holds unconditionally. For every original variety $X$ and every
prescribed proper closed subset $Z\subset X$, it produces an actual
alteration $X'\to X$ and an open immersion $X'\hookrightarrow\overline X'$
into a regular projective variety over $k$, with the actual image
of the inverse image of $Z$, together with the compactification
complement, equal to the support of an actual strict normal-crossings
divisor.
\end{theorem}
\begin{proof}
Apply the genuine whole original-field alteration descent theorem
to the preceding actual theorem over the actual algebraic closure
of $k$. That proved reduction chooses an actual dominating reduced
component, descends the actual finite-field model and entire
morphism, retains the same original boundary and all maps, and
descends the smooth ambient variety and strict normal-crossings
divisor to an actual finite field extension $E/k$. The resulting
ambient variety is regular and projective over the original $k$,
and the entire original prescribed inverse-image and compactification
complement has the native strict normal-crossings support.
It returns the same original native alteration statement over $k$.
No geometric integrality premise on the original variety is added.
\end{proof}

\endgroup

\begingroup

\begin{lemma}[Whole stalk ideals and their actual pullback]
\label{lem:node-fitting-ideal-sheaf-stalks}
\lean{DeJongAlterations.Algebra.NodeFittingIdealSheaf.exists_affineOpens_mem,DeJongAlterations.Algebra.NodeFittingIdealSheaf.stalkIdeal,DeJongAlterations.Algebra.NodeFittingIdealSheaf.map_germ_ideal_comap,DeJongAlterations.Algebra.NodeFittingIdealSheaf.map_germ_le_stalkIdeal,DeJongAlterations.Algebra.NodeFittingIdealSheaf.map_germ_eq_of_le,DeJongAlterations.Algebra.NodeFittingIdealSheaf.stalkIdeal_eq_map_germ,DeJongAlterations.Algebra.NodeFittingIdealSheaf.stalkIdeal_comap,DeJongAlterations.Algebra.NodeFittingIdealSheaf.ext_stalkIdeal}
\leanok
The supremum of the images of $I(U)$ under original germ maps over all
affine neighbourhoods $U$ of $x$ defines $I_x$. Every such single actual
$U$ computes this whole ideal. For every actual $g:Z\to X$,
$(g^*I)_z=I_{g(z)}\mathcal O_{Z,z}$ under the original stalk map.
Equal whole stalk ideals imply equal native ideal sheaf data; equal
inverse images on an actual open cover also imply equality.
\end{lemma}
\begin{proof}
For nested affine neighbourhoods, full-ideal restriction and composition
of the original germ maps give equal images. Any two have a common affine
basic neighbourhood, so all images in the supremum equal any one of them.
Apply the preceding affine pullback formula and the native germ--stalk-map
triangle. Native affine ideal extensionality detects equality on germs,
so equality on all stalks gives equality of whole chart ideals. An open
immersion has an isomorphism on actual stalks, whose ideal-extension map
is injective. Thus equal inverse images give equal stalk ideals on its
image, and the cover contains every point.
\end{proof}

\endgroup

\begingroup

\begin{lemma}[Kernel of the actual closed-subscheme stalk map]
\label{lem:node-fitting-ideal-sheaf-stalk-kernel}
\lean{DeJongAlterations.Algebra.NodeFittingIdealSheaf.ker_stalkMap_subschemeι}
\leanok
For any native ideal sheaf datum $I$ on any original scheme $X$, and any
actual $z\in V(I)$, the kernel of the original native stalk map
$\mathcal O_{X,\iota(z)}\to\mathcal O_{V(I),z}$ is exactly the whole
original stalk ideal $I_{\iota(z)}$.
\end{lemma}
\begin{proof}
Choose an actual affine neighbourhood $U$ of $\iota(z)$; its inverse
image in the native closed subscheme is affine. Both actual stalks are
the native localizations of those actual section rings at their actual
point primes. The original map on sections is surjective and has whole
kernel $I(U)$, by the native subscheme construction. The native
prime-ideal naturality identifies the actual source prime with the
contraction of the actual target prime. Surjectivity therefore makes the
image of the source prime complement exactly the target prime complement.
Native localization uniqueness and the actual germ--stalk-map triangle
identify the original stalk map with the native map of these actual
localizations. Its kernel is the image of the original whole section
kernel under the original source germ map, by the native localization
kernel theorem. The proved original stalk-ideal computation identifies
this image with the whole $I_{\iota(z)}$. No abstract replacement ring,
radical, Noetherian hypothesis or point-closedness is used.
\end{proof}

\endgroup

\begingroup

\begin{theorem}[Actual stalk quotient and original quotient-map triangle]
\label{thm:node-fitting-ideal-sheaf-actual-stalk-quotient}
\uses{lem:node-fitting-ideal-sheaf-stalk-kernel}
\lean{DeJongAlterations.Algebra.NodeFittingIdealSheaf.stalkQuotientEquiv,DeJongAlterations.Algebra.NodeFittingIdealSheaf.stalkQuotientEquiv_mk}
\leanok
The actual original closed-subscheme stalk is canonically ring-isomorphic
to $\mathcal O_{X,\iota(z)}/I_{\iota(z)}$, and the isomorphism sends every
original quotient class $[a]$ to the original native inclusion stalk map
applied to $a$.
\end{theorem}
\begin{proof}
The actual closed immersion is natively surjective on the original
stalks. Apply the native first isomorphism theorem to that same stalk
map, after identifying its whole actual kernel by the preceding lemma.
The quotient-by-equal-ideal map and the native quotient-kernel isomorphism
both preserve the original quotient representatives, giving exactly the
claimed full original map triangle.
\end{proof}

\endgroup

\begingroup

\begin{lemma}[Range of the fraction-chart inclusion]
\label{lem:affine-blowup-fraction-chart-range}
\uses{lem:affine-blowup-chart-over-base}
\lean{DeJongAlterations.affineBlowupFractionChartInclusion_opensRange}
\leanok
For $0\ne g\in I$, the native open immersion
$\operatorname{Spec}R[I/g]\to\operatorname{Bl}_I(\operatorname{Spec}R)$
has open range exactly $D_+(gT)$.
\end{lemma}
\begin{proof}
The chart inclusion is the spectrum of the homogeneous-chart ring
isomorphism followed by the native basic-open immersion. The first
map is an isomorphism, so precomposing by it leaves the range
unchanged. The range of the basic-open immersion is its defining
open $D_+(gT)$.
\end{proof}

\endgroup

\begingroup

\begin{lemma}[The local equation is a non-zero-divisor]
\label{lem:blowup-chart-equation-regular}
\uses{def:affine-blowup-fraction-chart}
\lean{DeJongAlterations.affineBlowupFractionChart_generator_isRegular}
\leanok
For any $g\ne0$, its image in $S$ is a non-zero-divisor.
\end{lemma}
\begin{proof}
The ring $S$ is a subring of the field $K$, hence a domain. The
composition $R\to S\to K$ is the injective fraction-field inclusion.
Consequently the image of $g$ in $S$ is nonzero. Multiplication by
a nonzero element of a domain is injective, on both sides, which is
the required regularity.
\end{proof}

\endgroup

\begingroup

\begin{lemma}[The exceptional ideal on the actual fraction chart]
\label{lem:blowup-exceptional-ideal-chart}
\uses{lem:affine-blowup-chart-over-base}
\lean{DeJongAlterations.affineBlowupExceptionalIdeal_comap_chart}
\leanok
Let $R$ be a domain and $0\ne g\in I$. On the actual chart
$\operatorname{Spec}R[I/g]$, the pulled-back blowup ideal is exactly the
ideal sheaf associated with $(g)$.
\end{lemma}
\begin{proof}
Ideal-sheaf pullback composes. The actual chart inclusion followed by
the blowup structure map is the spectrum map of $R\to R[I/g]$.
The affine ideal-sheaf pullback theorem identifies the result with the
extension of the original ideal, which is exactly $(g)$ by the checked
chart ideal computation.
\end{proof}

\endgroup

\begingroup

\begin{theorem}[Surjectivity of the actual graded map]\label{thm:rees-polynomial-surjective}
\lean{DeJongAlterations.Algebra.ReesPresentation.reesPolynomialMap_surjective}
\leanok
If the family $g$ generates $I$, the actual graded map is surjective.
\end{theorem}
\begin{proof}
Write an arbitrary Rees polynomial as its finite sum
$\sum_n a_nt^n$. Its membership in the Rees algebra says exactly that
$a_n\in I^n$ for every $n$. Apply the preceding lemma to each coefficient
in its finite support, obtaining homogeneous preimages $P_n$.
The finite sum $\sum_nP_n$ maps to the original polynomial by additivity.
Equality is equality in the actual Rees subalgebra of $R[t]$.
\end{proof}

\endgroup

\begingroup

\begin{definition}[Coefficient maps of native projective space]
\label{def:blowup-projective-coefficient-map}
\leanok
\lean{DeJongAlterations.Geometry.BlowupProjectiveEmbedding.polynomialCoefficientMap,DeJongAlterations.Geometry.BlowupProjectiveEmbedding.polynomialCoefficientMap_irrelevant,DeJongAlterations.Geometry.BlowupProjectiveEmbedding.projectiveCoefficientMap}
A ring homomorphism $R\to A$ induces the genuine coefficientwise graded
polynomial homomorphism and hence a native scheme morphism
$\mathbb P^m_A\to\mathbb P^m_R$.
\end{definition}
\begin{proof}
Coefficientwise mapping preserves every homogeneous degree. The target
irrelevant ideal is generated by its variables, each of which is the image
of the corresponding source variable in the source irrelevant ideal.
Thus the actual graded map satisfies the irrelevant-ideal containment
needed for native Proj. Applying native Proj gives the specified morphism.
\end{proof}

\endgroup

\begingroup

\begin{definition}[The affine map from the actual ideal generators]
\label{def:blowup-projective-affine-map}
\leanok
\lean{DeJongAlterations.Geometry.BlowupProjectiveEmbedding.generatorPolynomialMap,DeJongAlterations.Geometry.BlowupProjectiveEmbedding.generatorPolynomialMap_surjective,DeJongAlterations.Geometry.BlowupProjectiveEmbedding.affineGeneratorEmbedding,DeJongAlterations.Geometry.BlowupProjectiveEmbedding.affineGeneratorEmbedding_isClosedImmersion,DeJongAlterations.Geometry.BlowupProjectiveEmbedding.affineGeneratorEmbedding_structure,DeJongAlterations.Geometry.BlowupProjectiveEmbedding.affineProjectiveMap,DeJongAlterations.Geometry.BlowupProjectiveEmbedding.affineProjectiveMap_structure}
For a finite family $g_0,\ldots,g_m$ generating an actual ideal $I\subset A$,
the degree-one Rees elements $g_iT$ define an actual native map from
$\operatorname{Bl}_I(\operatorname{Spec}A)$ to $\mathbb P^m_R$.
\end{definition}
\begin{proof}
The actual graded polynomial evaluation $A[Y_0,\ldots,Y_m]\to\mathcal R(I)$
is surjective because the $g_i$ generate $I$. Its native Proj map is the
actual closed immersion into $\mathbb P^m_A$. Compose with the preceding
coefficient map $\mathbb P^m_A\to\mathbb P^m_R$. The structure square
follows from the two original native structure identities.
\end{proof}

\endgroup

\begingroup

\begin{lemma}[The complete pulled-back chart ideal]
\label{lem:blowup-projective-pulled-forms}
\leanok
\lean{DeJongAlterations.Geometry.BlowupProjectiveEmbedding.ideal_comap_affine,DeJongAlterations.Geometry.BlowupProjectiveEmbedding.coordinateAffineOpen,DeJongAlterations.Geometry.BlowupProjectiveEmbedding.coordinateIdeal_span,DeJongAlterations.Geometry.BlowupProjectiveEmbedding.pulledForm,DeJongAlterations.Geometry.BlowupProjectiveEmbedding.pulledForm_span}
For an actual affine open $U\subset j^{-1}D_+(X_i)$, the restrictions of
$F_a/X_i^{\deg F_a}$ generate the entire native pullback ideal on $U$.
\end{lemma}
\begin{proof}
Native pullback on arbitrary affine source and target opens corresponds to
extension of ideals along the actual map on sections: compose their canonical
spectrum maps with the original morphism, use the actual commutative spectrum
square, and apply the affine ideal-sheaf equivalence and its pullback formula.
On a projective coordinate chart, the exact homogeneous-quotient kernel is the
dehomogenized original ideal. The actual coordinate-to-section ring isomorphism
therefore carries the span of the dehomogenized forms onto the whole chart ideal.
Extending once more along the actual map on sections gives exactly the span of
the restricted forms and the original pulled-back native ideal, including all
nilpotent structure.
\end{proof}

\endgroup

\begingroup

\begin{lemma}[The actual subordinate blowup cover]
\label{lem:blowup-projective-subordinate-cover}
\leanok
\uses{lem:blowup-projective-pulled-forms}
\lean{DeJongAlterations.Geometry.BlowupProjectiveEmbedding.SubordinateAffineSite,DeJongAlterations.Geometry.BlowupProjectiveEmbedding.subordinateCover,DeJongAlterations.Geometry.BlowupProjectiveEmbedding.subordinateCover_locallyDirected}
The native affine Rees charts above affine opens subordinate to the
original projective coordinate opens form a locally directed open cover
of the actual global blowup.
\end{lemma}
\begin{proof}
For each point of the blowup, the image in projective space lies in a
coordinate open. The affine-open basis supplies an affine neighborhood
of its original base image contained in that inverse image. The exact
native range formula for the Rees chart includes the original blowup
point, proving coverage. Given a point of the intersection of two charts,
the full native locally directed affine-site cover supplies a smaller chart
mapping to both. Because it is contained in the first chart's base open,
it inherits that chart's projective-coordinate containment. Thus it lies
in the subordinate subcategory, and all transition identities are the
original native affine Rees transition identities.
\end{proof}

\endgroup

\begingroup

\begin{theorem}[The projective morphism from the native global blowup]
\label{thm:blowup-projective-global-map}
\source{dejong96-alterations:page-0018}
\leanok
\lean{DeJongAlterations.Geometry.BlowupProjectiveEmbedding.subordinateCoordinate,DeJongAlterations.Geometry.BlowupProjectiveEmbedding.subordinateCoordinate_le,DeJongAlterations.Geometry.BlowupProjectiveEmbedding.formsChartMap,DeJongAlterations.Geometry.BlowupProjectiveEmbedding.formsChartMap_compatible,DeJongAlterations.Geometry.BlowupProjectiveEmbedding.projectiveBlowupMapOfForms,DeJongAlterations.Geometry.BlowupProjectiveEmbedding.projectiveBlowupMapOfForms_restrict,DeJongAlterations.Geometry.BlowupProjectiveEmbedding.projectiveBlowupMapOfForms_structure,DeJongAlterations.Geometry.BlowupProjectiveEmbedding.exists_projectiveBlowupMap}
\uses{lem:blowup-projective-subordinate-cover}
Let $j:X\hookrightarrow\mathbb P^n_R$ be an actual projective embedding and
let an original ideal $D$ be represented by finite forms $F_a$ of one degree
$d$. The genuine local generator maps glue to a native morphism
$\operatorname{Bl}_D X\to\mathbb P^m_R$ over $R$.
In particular, over a Noetherian $R$, the actual projective embedding and
$D$ alone internally produce such a finite $m$ and morphism, without a
supplied homogeneous representation.
\end{theorem}
\begin{proof}
On the affine inverse image of $D_+(X_i)$, use the actual functions
$F_a/X_i^d$ generating the whole original ideal. Construct the preceding
affine projective map on its actual Rees blowup. On an overlap with the
chart selected by $X_j$, each generator is multiplied by the same unit
$(X_i/X_j)^d$. The preceding lemma gives equality of the actual scheme
morphisms after restriction. Native scheme gluing on the genuine blowup
cover supplies the global morphism, and the local structure identities
glue to its original base identity.
\end{proof}

\endgroup

\begingroup

\begin{lemma}[Native closed immersion on each actual target chart]
\label{lem:segre-native-chart-closed-immersion}
\lean{DeJongAlterations.Geometry.Segre.coordinateChartIso,DeJongAlterations.Geometry.Segre.coordinateChartInclusion_structure,DeJongAlterations.Geometry.Segre.chartClosedImmersion,DeJongAlterations.Geometry.Segre.chartClosedImmersion_isClosedImmersion,DeJongAlterations.Geometry.Segre.chartToProjective,DeJongAlterations.Geometry.Segre.chartToProjective_structure}
\leanok
The actual affine scheme
$\operatorname{Spec}(A_i\otimes_R B_j)$ has a closed
immersion into the native chart $D_+(z_{ij})\subset
\mathbf P_R^N$, induced by the displayed coordinate map.
Its structure map agrees with the original $R$-coefficient map.
\end{lemma}
\begin{proof}
The checked native projective chart isomorphism identifies
$D_+(z_{ij})$ with $\operatorname{Spec}C_{ij}$.
The actual surjective ring map of the preceding lemma
therefore induces the claimed closed immersion.
All algebra maps are $R$-algebra maps, so their coefficient
triangles commute. The native Proj chart structure
triangle identifies that coefficient triangle with the
restriction of $\mathbf P_R^N\to\operatorname{Spec}R$.
\end{proof}

\endgroup

\begingroup

\begin{lemma}[Actual native coordinate overlap transitions]
\label{lem:segre-native-coordinate-overlap}
\uses{lem:segre-native-chart-closed-immersion}
\lean{DeJongAlterations.Algebra.Segre.CoordinateOverlapRing,DeJongAlterations.Algebra.Segre.overlapCoefficients,DeJongAlterations.Algebra.Segre.overlapAlgebra,DeJongAlterations.Algebra.Segre.overlapLeft,DeJongAlterations.Algebra.Segre.overlapRight,DeJongAlterations.Algebra.Segre.homogeneousChartRatioHom_coordinate,DeJongAlterations.Algebra.Segre.overlap_coordinate_cross,DeJongAlterations.Algebra.Segre.overlapLeft_pivot_isUnit}
\leanok
In the actual homogeneous localization at $x_i x_k$, the two
constructed coordinate-chart restriction maps satisfy
\[
 (x_\ell/x_k)(x_k/x_i)=x_\ell/x_i.
\]
In particular the change-of-pivot coordinate $x_k/x_i$ is a unit
on this actual overlap, with inverse $x_i/x_k$.
\end{lemma}
\begin{proof}
The native homogeneous localization maps send $x_\ell/x_i$ to
$(x_\ell x_k)/(x_i x_k)$ and $x_\ell/x_k$ to
$(x_\ell x_i)/(x_i x_k)$. Their values embed injectively into the
ordinary localization at $x_i x_k$. Cross-multiplication there proves
the displayed identity on actual whole-ring elements. Set $\ell=i$
to obtain the stated actual inverse. All maps preserve the original
coefficient map, including over rings with nilpotents.
\end{proof}

\endgroup
